\documentclass[10pt,a4paper]{amsart}
\usepackage[margin=19mm]{geometry}
\usepackage{amsmath,amssymb,mathtools}
\usepackage[scr=rsfs,cal=euler]{mathalfa}
\usepackage{xfrac}
\usepackage[unicode,hidelinks,bookmarks=false]{hyperref}
\newcommand{\ie}{i.e.\ }
\newcommand{\cf}{cf.\ }
\newcommand{\resp}{respectively\ }
\newcommand{\Subsection}{\subsection}
\providecommand{\nobreakdash}{\nobreak\hbox{-}}
\setlength{\emergencystretch}{2em}
\multlinegap0pt
\usepackage{amssymb}
\newtheorem{lemma}{Lemma}
\title{A new simple family of non-periodic tilings with square tiles}
\author{Nikolay Vereshchagin}
\date{Source: arXiv:2307.16134v3 (CC BY 4.0)}
\begin{document}
\maketitle

\noindent\emph{Source and licence.} A new simple family of non-periodic tilings with square tiles. Version arXiv:2307.16134v3 (CC BY 4.0), distributed by arXiv under the Creative Commons Attribution 4.0 International licence, http://creativecommons.org/licenses/by/4.0/, as recorded in the arXiv licence metadata for this exact version and shown on the abstract page https://arxiv.org/abs/2307.16134v3. Full licence notice: \texttt{licenses/arxiv-2307.16134-CC-BY-4.0.txt}.\par\smallskip

\noindent\textit{Editorial scope.} Counted unit: the article's lemma lemma (source0.tex lines 350--352). The article's complete original proof of it is delivered with it (source0.tex lines 353--422, one \texttt{proof} environment of 70 lines). No mathematics is written, repaired or re-derived: the statement and the proof are the article's own lines, in the article's own notation, and no retained line is substituted or re-flowed.  No further source-local context is needed: every cross-reference of every retained line resolves inside the two delivered spans.  One declared adaptation: exactly 4 ancillary strings inside the retained spans are omitted (image-inclusion commands and, where the source repeats a label, the repeated label definitions) and no other line differs from the source; the figure environments, with their placement, captions and labels, are kept, and the package records the omitted strings in fidelity receipts bound to the affected bodies.  No retained line contains a citation, so the wrapper carries no bibliography and no bibliography entry.  The retained lines contain the article's own diagram code, which the wrapper typesets with the standard tikz packages; no image file is delivered and the package declares no asset dependency.  Editorial formatting only, in addition: an AMS \texttt{amsart} preamble that restates the article's own declarations and macros used by the retained lines, namely the environments \texttt{lemma} on separate counters, together with the standard packages those lines need.  The retained environments print in this document as Lemma 1 in order of appearance, whereas the article prints the same items with the numbering of its own full document.  The complete native source of the article as distributed in the arXiv e-print for this exact version is bundled unchanged as \texttt{sources/144956/source0.tex}.
\begin{lemma}
Any supertile is a correct tiling.
\end{lemma}

\begin{proof}
This is proved by induction. The induction base is trivial, since any tile constitutes a correct tiling.

For the inductive step, we have to show that the decomposition of a correct tiling $T$ is again correct. 
But there is a problem here. This is true for tilings of the plane, but not for tilings of its parts. 
For example, a tiling consisting of two triangles with a common leg in Fig.~\ref{pic27},\begin{figure}[ht]
\begin{center}

\end{center}
\caption{The decomposition a correct tiling of a part of the plane may not be correct.}\label{pic27}
\end{figure}
is correct (the third condition is satisfied because there are no internal vertices), but its decomposition is not. 
To make the inductive step, it is necessary to impose a restriction on the crossings at the boundary points. 
To prove the lemma, it will suffice to require that the crossings with the center at any vertex of the boundary, 
except for the extreme ones, are halves of the legal crossings, that is, they have one of the two forms in Fig.~\ref{pic26}.
\begin{figure}[ht]
\begin{center}

\vskip .5cm

\end{center}
\caption{Allowed crossings on the boundary of a supertile. The  purple sides must have the same color, red or green.
}\label{pic26}
\end{figure}

Now we can make the inductive step. 
Assume that $T$ is a supertile and is correct. We have to show that so is its decomposition $\sigma T$.
It is obvious that the first condition 
(tiles border side by side and colors and orientations of shared sides match) %for  $\sigma T$
is inherited during decomposition.

Let us verify the second condition (adjacent triangles have different colors) for  $\sigma T$. 
Let $F$ be a tile in $\sigma T$ 
and let $G$ be its sibling; both arise from decomposing a single tile $I$ of $T$. W.l.o.g. assume that  $F$ is a right, and hence green, tile (Fig.~\ref{pic28}).
\begin{figure}[ht]
\begin{center}

\end{center}
\caption{All neighbors of a green tile $F$ in the decomposition are red.}\label{pic28}
\end{figure}
We need to prove that all neighbors of $F$ are red.

The tile $G$ is red by definition of the substitution.
Besides $G$, the tile $F$ can have other neighbors.
Namely, if the hypotenuse of $I$ does not lie on the edge of $T$, then the tiling $T$ contains the tile $E$ bordering $I$ along the hypotenuse.
As a result of cutting $E$, two tiles will appear, of which it is the left, and therefore red, tile that borders $F$ along the leg.

Moreover, if the leg of the tile $I$ does not lie on the boundary of the supertile, then the tiling $T$ contains the tile $H$,
which is the reflection of the tile $I$ with respect to this leg.
%The tile $H$ is oriented exactly as it is drawn,
%that is, its right angle adjoins the vertex $A$.
Indeed, otherwise the crossing centered at the vertex $A$ (Fig.~\ref{pic28}) would be illegal
(right and acute angles cannot converge at a legal crossing).
When the tile $H$ is decomposed, it is the left, and therefore red, tile that is adjacent to $F$.

Let us check the third condition for the tiling $\sigma T$ (that all crossings are legal).
To do this, we need to check that during decomposition, a legal crossing is converted into a legal one,
and the same holds for half crossings in Fig.~\ref{pic26}.
This is done by a direct verification.
In addition, we need to check that all the crossings at the new vertices are legal as well.

New vertices are the centers of  hypotenuses of tiles from $T$.
The hypotenuse itself becomes the axis of the new crossing, and the new crossing is $C_4$ by the definition of the substitution.

 
Finally, we check the fourth condition. The crossings of the form $C_4$ are only the crossings at the centers of the hypotenuses of the tiles in $T$,
and by the definition of substitution, all the left triangles in them are red and the right ones are green.

The third and fourth conditions for crossings centered on the boundary are verified in a similar way.
\end{proof}

\end{document}
